% !TeX root = ../exam-zh-doc-basic.tex

\section{常见场景}

本章介绍实际使用中的常见需求和解决方案。

\subsection{答案控制方案}

\subsubsection{生成学生版和教师版}

最常见的需求是生成两个版本的试卷：

\textbf{学生版（默认）：}

什么都不设置，默认就是学生版，所有答案都隐藏。

\textbf{教师版：}

在文档开头添加：

\begin{latexcode}
  \examsetup{
    question/show-answer = true,        % 显示选择题和填空题答案
    solution/show-solution = show-stay, % 显示解答题答案
    question/show-points = true         % 显示分值（可选）
  }
\end{latexcode}

\begin{tcolorbox}[colback=green!5!white, colframe=green!75!black, title=\textbf{最佳实践：版本切换方法对比}]
\textbf{方法一：注释法（推荐新手）}

适合场景：偶尔切换版本，手动编译

\begin{latexcode}
% 生成学生版：注释掉下面的代码
% 生成教师版：取消注释
% \examsetup{
%   question/show-answer = true,
%   solution/show-solution = show-stay
% }
\end{latexcode}

\textbf{方法二：条件编译（推荐进阶用户）}

适合场景：频繁切换版本，或需要自动化编译

\begin{latexcode}
% 在文档开头定义开关
\newif\ifteacher
\teachertrue   % 教师版：改为 \teachertrue
% \teacherfalse  % 学生版：改为 \teacherfalse

\ifteacher
  \examsetup{
    question/show-answer = true,
    solution/show-solution = show-stay
  }
\fi
\end{latexcode}

\textbf{方法三：命令行参数（推荐批量生成）}

适合场景：自动化脚本，一次编译生成两个版本

创建主文件 \file{exam-main.tex}：
\begin{latexcode}
\ifdefined\SHOWANSWER
  \examsetup{
    question/show-answer = true,
    solution/show-solution = show-stay
  }
\fi
\end{latexcode}

命令行编译：
\begin{shellcode}
# 学生版
xelatex exam-main.tex

# 教师版
xelatex "\def\SHOWANSWER{1}\input{exam-main.tex}"
\end{shellcode}
\end{tcolorbox}

\subsubsection{解答显示位置}

解答题的答案可以在原位显示，也可以移到试卷最后：

\begin{latexcode}
  \examsetup{
    % solution/show-solution = hide,       % 隐藏（学生版）
    % solution/show-solution = show-stay,  % 原位显示（教师版）
    solution/show-solution = show-move   % 移到最后
  }
\end{latexcode}

|show-move| 模式适合制作"答案另附"的试卷。


\subsection{页面设置}

\subsubsection{纸张大小}

\begin{latexcode}
  \examsetup{
    page/size = a4paper,  % A4 纸（默认）
    % page/size = a3paper,  % A3 纸
  }
\end{latexcode}

\subsubsection{页边距调整}

`\cls{exam-zh}` 会在导言区末尾设置默认页面参数，所以自定义页边距时建议把 |\geometry| 放进 |\AtEndPreamble|：

\begin{latexcode}[deletetexcs={\documentclass},morekeywords={\documentclass}]
  \documentclass{exam-zh}

  \AtEndPreamble{
    \geometry{
      top = 2.5cm,
      bottom = 2cm,
      left = 2cm,
      right = 2cm
    }
  }
\end{latexcode}

\subsubsection{分栏设置}

让试卷内容分成两栏：

\begin{latexcode}
  \usepackage{multicol}

  \begin{document}
    \title{标题}
    \maketitle

    \begin{multicols}{2}  % 开始两栏
      % 题目内容...
    \end{multicols}       % 结束两栏
  \end{document}
\end{latexcode}


\subsection{图文混排详解}

本节介绍 \cls{exam-zh} 自带的图文排版命令。

\subsubsection{基本用法}

\textbf{使用 \cmd{\textfigure}：}

\begin{latexcode}
  \textfigure{
    如图所示，已知三角形 ABC 的三边长分别为 3、4、5，
    求三角形的面积。
  }{
    \includegraphics[width=3cm]{triangle.pdf}
  }
\end{latexcode}

\textbf{说明：}
\begin{itemize}
  \item 第一个大括号是文字内容
  \item 第二个大括号是图片内容
  \item 更细的控制参数可以在完整文档中查询 |textfigure/*| 键值
\end{itemize}

\subsubsection{多张图片}

如果一题对应多张图片，可以使用 |multifigures| 环境：

\begin{latexcode}
  \begin{multifigures}[columns = 2]
    \item[题 1 图] \includegraphics[width=3cm]{fig1.png}
    \item[题 2 图] \includegraphics[width=3cm]{fig2.png}
  \end{multifigures}
\end{latexcode}

\textbf{适用场景：}
\begin{itemize}
  \item 一段文字配一张图：优先用 |\textfigure|
  \item 多张题图并排展示：优先用 |multifigures|
  \item 只想简单插图：直接用 |\includegraphics|
\end{itemize}

\subsubsection{简单插图}

\textbf{居中显示：}

\begin{latexcode}
  \begin{center}
    \includegraphics[width=0.5\textwidth]{fig.png}
  \end{center}
\end{latexcode}

更多示例见完整文档的“图文排版”章节。


\subsection{行距和间距调整}

本节介绍如何调整试卷的行距和各种间距，以适应不同需求（如节约纸张）。

\subsubsection{全局行距调整}

\textbf{方法 1：使用 \textbackslash linespread}

\begin{latexcode}[deletetexcs={\documentclass},morekeywords={\documentclass}]
  \documentclass{exam-zh}

  \linespread{0.9}  % 行距缩小到 90%（默认 1.0）

  \begin{document}
    % 试卷内容...
  \end{document}
\end{latexcode}

\textbf{方法 2：使用 setspace 宏包（更精细）}

\begin{latexcode}[deletetexcs={\documentclass},morekeywords={\documentclass}]
  \documentclass{exam-zh}

  \usepackage{setspace}
  \setstretch{0.85}  % 0.85 倍行距

  % 或者使用预定义命令
  % \singlespacing    % 单倍行距
  % \onehalfspacing   % 1.5 倍行距
  % \doublespacing    % 双倍行距
\end{latexcode}

\textbf{注意：}
\begin{itemize}
  \item 行距不要低于 0.8，否则阅读困难
  \item 建议范围：0.85--1.2
  \item 行距过小可能导致公式重叠
\end{itemize}

\subsubsection{题目间距调整}

\textbf{调整题目之间的垂直间距：}

\begin{latexcode}
  \examsetup{
    question/top-sep = 0.2em,
    question/bottom-sep = 0.2em,
  }
\end{latexcode}

\textbf{调整小题（problem）间距：}

\begin{latexcode}
  \examsetup{
    problem/top-sep = 0.5em,
    problem/bottom-sep = 0.5em,
  }
\end{latexcode}

\subsubsection{选项间距调整}

\textbf{调整选择题选项的垂直间距：}

\begin{latexcode}
  \examsetup{
    choices/linesep = 0.3em,  % 多行选项之间的额外间距
  }
\end{latexcode}

\textbf{调整选项的水平间距：}

\begin{latexcode}
  \examsetup{
    choices/column-sep = 1em,  % 列与列之间的最小间距
  }
\end{latexcode}

\subsubsection{段落间距调整}

\textbf{调整段落间距：}

\begin{latexcode}[deletetexcs={\documentclass},morekeywords={\documentclass}]
  \documentclass{exam-zh}

  % 减小段落间距
  \setlength{\parskip}{0.3em}

  % 取消首行缩进（可选）
  \setlength{\parindent}{0pt}
\end{latexcode}

\subsubsection{紧凑模式示例}

\textbf{综合示例：让试卷更紧凑以节约纸张}

\begin{latexcode}[deletetexcs={\documentclass},morekeywords={\documentclass}]
  \documentclass{exam-zh}

  % === 全局行距 ===
  \usepackage{setspace}
  \setstretch{0.9}

  % === 页边距 ===
  \usepackage{geometry}
  \geometry{
    top = 2cm,
    bottom = 1.5cm,
    left = 1.8cm,
    right = 1.8cm
  }

  % === 题目和选项间距 ===
  \examsetup{
    question/top-sep = 0.2em,
    question/bottom-sep = 0.2em,
    choices/linesep = 0.3em,
    choices/column-sep = 0.8em,
  }

  % === 段落间距 ===
  \setlength{\parskip}{0.2em}

  \begin{document}
    % 试卷内容...
  \end{document}
\end{latexcode}

\textbf{效果：}
\begin{itemize}
  \item 可节约约 20--30\% 的纸张
  \item 保持基本可读性
  \item 适合练习题或非正式考试
\end{itemize}

\textbf{注意事项：}
\begin{itemize}
  \item 正式考试建议保持默认设置
  \item 紧凑模式可能影响阅读体验
  \item 图片密集的试卷效果不明显
\end{itemize}

\subsubsection{局部调整间距}

如果只想在某些地方调整间距：

\begin{latexcode}
  % 增加间距
  \vspace{1em}

  % 减少间距
  \vspace{-0.5em}

  % 强制分页前的弹性间距
  \vfill
\end{latexcode}

更多排版技巧见完整文档。


\subsection{特殊功能}

\subsubsection{草稿纸}

在试卷最后添加草稿纸：

\begin{latexcode}
  \newpage
  \section*{草稿纸}
  \draftpaper
\end{latexcode}

\subsubsection{方格纸}

需要一串可填写的小方格时，可以使用 |\examsquare|：

\begin{latexcode}
  考号：\examsquare{10}
\end{latexcode}

\subsubsection{评分框}

在试卷中添加评分表：

\begin{latexcode}
  \scoringbox
\end{latexcode}

详细配置见完整文档。

\subsubsection{密封线}

隐藏密封线：

\begin{latexcode}
  \examsetup{
    sealline/show = false
  }
\end{latexcode}

修改密封线文字：

\begin{latexcode}
  \examsetup{
    sealline/text = {密封线内不要答题}
  }
\end{latexcode}

更多密封线设置见完整文档。


\subsection{学科专用功能}

\subsubsection{语文试卷}

\textbf{古诗排版：}

使用 |poem| 环境：

\begin{latexcode}
  \begin{poem}
    白日依山尽，\\
    黄河入海流。\\
    欲穷千里目，\\
    更上一层楼。
  \end{poem}
\end{latexcode}

\textbf{阅读材料：}

使用 |material| 环境：

\begin{latexcode}
  \begin{material}[title = {文章标题}, author = {作者}]
    文章内容...
  \end{material}
\end{latexcode}

\subsubsection{英语试卷}

\textbf{作文答题框：}

\begin{latexcode}
  \begin{writingbox}[title = {My Dream}]
    \vspace*{10em}
  \end{writingbox}
\end{latexcode}


\subsection{快速参考}

\textbf{常用配置速查：}

\begin{latexcode}
  \examsetup{
    % === 答案控制 ===
    question/show-answer = true,  % 显示选择题和填空题答案
    solution/show-solution = show-stay,  % 显示解答题答案

    % === 分值显示 ===
    question/show-points = true,  % 显示题目分值

    % === 选项排列 ===
    choices/columns = 2,  % 选项排成2列

    % === 页面设置 ===
    sealline/show = false,  % 隐藏密封线
  }
\end{latexcode}

更多高级功能和详细参数，请查阅完整文档 \file{exam-zh-doc.pdf}。
